\chapter{Reference manual: orientation}
\label{ch:refman}

\begin{aside}
The rest of this book taught you to think in \maya{}. This
appendix is the dictionary you reach for when you've already
internalised the grammar and just need to remember the exact
spelling of a verb.
\end{aside}

Every entry that follows obeys the same shape:

\begin{itemize}[leftmargin=*]
  \item A one-line gloss of what the thing is.
  \item The signature you'll type.
  \item A small but \emph{complete} program. Not a snippet —
    a file you can save, build, and run. Every example
    starts with the includes and ends with a \code{main()}
    or equivalent.
  \item A short \emph{notes} section for the rough edges.
\end{itemize}

\section{How to read these examples}

Every program in this manual compiles against the same
single header:

\begin{lstlisting}[language=C++]
#include <maya/maya.hpp>
using namespace maya;
using namespace maya::dsl;
\end{lstlisting}

Some widgets pull an extra header (\code{<maya/widget/...>}).
Those are listed inline where they appear.

The build line is always the same:

\begin{lstlisting}[language=bash]
g++ -std=c++23 example.cpp -lmaya -o example
\end{lstlisting}

\noindent or, with CMake:

\begin{lstlisting}[language=cmake]
find_package(maya REQUIRED)
add_executable(example example.cpp)
target_link_libraries(example PRIVATE maya::maya)
\end{lstlisting}

\section{Layout of this appendix}

Twenty-eight chapters. Each is a slice of the surface area:

\begin{itemize}[leftmargin=*]
  \item Programs and the runtime loop.
  \item Elements: text, layout, decoration.
  \item Style: colour, weight, modifiers.
  \item Subscriptions: keys, timers, signals.
  \item Commands: side effects as data.
  \item Widgets: every shipped widget, with a working demo.
  \item Inline rendering: \code{print} and \code{live}.
  \item Edge cases: resize, exit, errors, focus.
\end{itemize}

\section{The example template}

Use this as a starting point when you copy from this manual:

\begin{lstlisting}[language=C++]
#include <maya/maya.hpp>

using namespace maya;
using namespace maya::dsl;

struct App {
    struct Model {};
    struct Quit {};
    using Msg = std::variant<Quit>;

    static Model init() { return {}; }

    static auto update(Model m, Msg)
        -> std::pair<Model, Cmd<Msg>>
    {
        return {m, Cmd<Msg>::quit()};
    }

    static Element view(const Model&) {
        return t<"hello"> | Bold;
    }

    static auto subscribe(const Model&) -> Sub<Msg> {
        return key_map<Msg>({ {'q', Quit{}} });
    }
};

int main() { run<App>({.title = "demo"}); }
\end{lstlisting}

\noindent Every working example in this appendix is a
variation on that skeleton. If you forget where to start,
start here.

\section{A note on the prose}

These chapters are deliberately terser than the rest of the
book. You're here for the API, not the philosophy. Bullets
where bullets work; code where code works; one paragraph
of context per chapter, then on to the next entry.

\section{Recap}

\begin{recap}
\begin{itemize}[leftmargin=*]
  \item Every example is a complete program.
  \item Every program follows the same skeleton.
  \item Twenty-eight chapters cover every public surface.
  \item Read on demand, not in order.
\end{itemize}
\end{recap}

\section*{Workshop}

\begin{enumerate}[leftmargin=*]
  \item Type the template above into a file and build it.
  \item Run it; press \code{q}.
  \item You now have a working \maya{} project. The rest of
    this appendix is just things you can substitute in.
\end{enumerate}
